\documentclass[10pt,a4paper]{amsart}
\usepackage[margin=19mm]{geometry}
\usepackage{amsmath,amssymb,mathtools}
\usepackage[scr=rsfs,cal=euler]{mathalfa}
\usepackage{xfrac}
\usepackage[unicode,hidelinks,bookmarks=false]{hyperref}
\newcommand{\ie}{i.e.\ }
\newcommand{\cf}{cf.\ }
\newcommand{\resp}{respectively\ }
\newcommand{\Subsection}{\subsection}
\providecommand{\nobreakdash}{\nobreak\hbox{-}}
\setlength{\emergencystretch}{2em}
\multlinegap0pt
\usepackage{float}
\usepackage{amssymb}
\usepackage{dsfont}
\usepackage{tikz}
\usepackage{xcolor}
\newtheorem{definition}{Definition}
\newtheorem{lemma}{Lemma}
\title{Moving Least Squares without Quasi-Uniformity: A Stochastic Approach}
\author{Shir Tapiro Moshe \and Yariv Aizenbud \and Barak Sober}
\date{Source: arXiv:2601.13782v5 (CC BY 4.0)}
\begin{document}
\maketitle

\noindent\emph{Source and licence.} Moving Least Squares without Quasi-Uniformity: A Stochastic Approach. Version arXiv:2601.13782v5 (CC BY 4.0), distributed by arXiv under the Creative Commons Attribution 4.0 International licence, http://creativecommons.org/licenses/by/4.0/, as recorded in the arXiv licence metadata for this exact version and shown on the abstract page https://arxiv.org/abs/2601.13782v5. Full licence notice: \texttt{licenses/arxiv-2601.13782-CC-BY-4.0.txt}.\par\smallskip

\noindent\textit{Editorial scope.} Counted unit: the article's lemma lambda\_min (source0.tex lines 862--869). The article's complete original proof of it is delivered with it (source0.tex lines 871--980, one \texttt{proof} environment of 110 lines). No mathematics is written, repaired or re-derived: the statement and the proof are the article's own lines, in the article's own notation, and no retained line is substituted or re-flowed.  Retained uncounted as the source-local context the proof needs: lines 240--244; lines 472--475; lines 757--760; lines 792--795.  Nothing was omitted from any retained span, so the package carries no omission record.  No retained line contains a citation, so the wrapper carries no bibliography and no bibliography entry.  No retained line contains a figure environment, an image-inclusion command or drawing code, so no image is delivered and the package declares no asset dependency.  Editorial formatting only, in addition: an AMS \texttt{amsart} preamble that restates the article's own declarations and macros used by the retained lines, namely the environments \texttt{definition}, \texttt{lemma} on separate counters, together with the standard packages those lines need.  The retained environments print in this document as Definition 1, Lemma 1 in order of appearance, whereas the article prints the same items with the numbering of its own full document.  The complete native source of the article as distributed in the arXiv e-print for this exact version is bundled unchanged as \texttt{sources/193044/source0.tex}.
\begin{definition} ("nicely behaving distribution").
    Let $\mu$ denote the uniform measure over $\Omega$, induced by the Lebesgue measure.
A probability density function $p(x)$ on $\Omega$ is called nicely behaving distribution, if there exist constants $0<\underline{c} \leq \bar{c}<\infty$ such that $\underline{c} \cdot \mu \leq p(x) \leq \bar{c} \cdot \mu$.
    \label{NicelyDef}
\end{definition}

\begin{equation}
    \mathscr{E}_I:=\{ \# I_{B(\hat{x}, s h_n)} \in [\gamma_1 \log n,\gamma_2 \log n] \},
\label{E_I event}
\end{equation}

\begin{equation}
    \mathscr{A}_{n \alpha \beta}(x) =\frac{1}{N_{B\left(x, h_n\right)}} \sum_{I_{B\left(x, h_n\right)}} \theta_h(x_i-x)\frac{\left(x_i-x\right)^\alpha\left(x_i-x\right)^\beta}{h_n^{|\alpha|+|\beta|}}. 
    \label{Anab}
\end{equation}

\begin{equation}
    P\left(\left|\mathscr{A}_{n a \beta}^w(x)-\mathscr{A}^w_{\alpha \beta}\right| \leq \varepsilon \Big|  \mathscr{E}_I\right)\geq 1-2 e^{-c'\cdot \varepsilon^2 \log(n) }=1-\frac{2}{n^{c'\cdot \varepsilon^2 }},
    \label{eq:bernstein inequality}
\end{equation}

\begin{lemma} 
Given the event $\mathscr{E}_I$ defined in \eqref{E_I event} holds, there exists a constant $C_\lambda>0$,  such that for a fixed but arbitrary evaluation point $x\in \Omega$,
$$
P\left(\lambda_{\min }\left(\mathscr{A}_n(x)\right) \geq C_\lambda \Big| \mathscr{E}_I\right)\geq 1-\frac{2A^2}{n^\phi},
$$
where $\phi$ is a real positive constant, $A$ is the dimension of the matrix $\mathscr{A}_n(x)$ and $\lambda_{\text {min }}\left(\mathscr{A}_n(x)\right)$ is the smallest eigenvalue of $\mathscr{A}_n(x)$. In particular, with the stated probability $\mathscr{A}_n(x)$ is positive definite and full rank.
    \label{lemma 3.7 davoud - lambda_min}
\end{lemma}

\begin{proof} 
Let $\pi=\left(\pi_\alpha\right)_{\alpha \in A}$ be a unit vector in $\mathbb{R}^{|A|}$ such that $\sum_{\alpha\in A} \pi_\alpha^2=1$. For $z\in\mathbb{R}^d$, we define the polynomial:
$$
P(z)=\sum_{\alpha \in A} \pi_\alpha z^\alpha
$$
Now define the empirical quadratic form:
$$
\mathcal{Q}_n(\pi):=\pi^{\top} \mathscr{A}_n(x) \pi.
$$
From \eqref{Anab} we get
$$
\mathcal{Q}_n(\pi)=\frac{1}{N_{B\left(x, h_n\right)}} \sum_{i \in I_{B\left(x, h_n\right)}} \theta_h\left(x_i-x\right) \cdot\left(\sum_\alpha \pi_\alpha \frac{\left(x_i-x\right)^\alpha}{h_n^{|\alpha|}}\right)^2 .
$$
Let us define the rescaled coordinates:
$$
t_i:=\frac{x_i-x}{h_n}, \quad \text { so that } \quad x_i=x+h_n t_i.
$$
Substituting $t_i=\frac{x_i-x}{h_n}$, and using $\theta_h\left(x_i-x\right)=\Phi\left(t_i\right)$, we get:

$$
\mathcal{Q}_n(\pi)=\frac{1}{N_{B\left(x, h_n\right)}} \sum_{i \in I_{B\left(x, h_n\right)}} \Phi\left(t_i\right) \cdot\left(\sum_\alpha \pi_\alpha t_i^\alpha\right)^2=\frac{1}{N_{B\left(x, h_n\right)}} \sum_{i \in I_{B\left(x, h_n\right)}} \Phi\left(t_i\right) \cdot P\left(t_i\right)^2,
$$
where $P(t)=\sum_\alpha \pi_\alpha t^\alpha$.
Given the event $\mathscr{E}_I$, the number of samples $\left\{t_i\right\}_{i \in I_{B\left(x, h_n\right)}}$ inside the rescaled ball $B(0, s)$ goes to infinity  as $n \rightarrow \infty$ in probability.  We now show that 
$\mathbb{E}\left[\mathcal{Q}_n(\pi)\right]$ is positive.
$$
\mathbb{E}\left(\mathcal{Q}_n(\pi)\right)=\frac{\int_{B\left(x, h_n\right)} \Phi\left(\frac{u-x}{h_n}\right) P\left(\frac{u-x}{h_n}\right)^2 p(u) d u}{\int_{B\left(x, h_n\right)} p(u) d u}.
$$
By applying the change of variables $t=\frac{u-x}{h_n}$, which means $d u=h_n^d d t$, we obtain:
$$
\mathbb{E}(\mathcal{Q}_n(\pi))=\frac{\int_{B(0, s)} \Phi(t) P(t)^2 p\left(x+h_n t\right) h_n^d d t}{\int_{B(0, s)} p\left(x+h_n t\right) h_n^d d t}.
$$
Recall now that $\int_{B(0,s)} \underline{c} d \mu \leq \int_{B(0,s)} d p \leq\int_{B(0,s)} \bar{c} d \mu$, where $p$ is a density function satisfies Definition \ref{NicelyDef}, hence,
\begin{equation}
\begin{aligned}
    \underline{\beta} \int_{B(0, s)} \Phi(t) P(t)^2 dt 
    &= \frac{\underline{c}\int_{B(0, s)} \Phi(t) P(t)^2 d t}{\bar{c}\int_{B(0, s)} d t} 
    \leq \mathbb{E}(\mathcal{Q}_n(\pi)) \\
    &\leq \frac{\bar{c}\int_{B(0, s)} \Phi(t) P(t)^2 d t}{\underline{c}\int_{B(0, s)} d t} 
    = \bar{\beta} \int_{B(0, s)} \Phi(t) P(t)^2 dt.
\end{aligned}
\end{equation}
where, $\underline{\beta}= \frac{\underline{c}}{\bar{c}\cdot \text{Vol}B(0,s)} $ and $\bar{\beta}=\frac{\bar{c}}{\underline{c}\text{Vol}B(0,s)} $. Denote,
$$
\mathcal{Q}(\pi)=\int_{B(0, s)} \Phi(t) \cdot P(t)^2 dt.
$$
Thus,
\begin{equation}
    \underline{\beta} \mathcal{Q}(\pi) \leq \mathbb{E}(\mathcal{Q}_n(\pi))\leq
\bar{\beta} \mathcal{Q}(\pi).
\label{bound of EQ_n}
\end{equation}
Since the integrand $\Phi(t) \cdot P(t)^2 $ is non-negative, and $P(t) $ is not identically zero on $ B(0,s)$, the integral $\mathcal{Q}(\pi) $ is strictly positive for all $\|\pi\|=1 $, i.e., $\mathcal{Q}(\pi)>0$.
Moreover, $\mathcal{Q}(\pi)$ is continuous in $\pi$, and is defined on a compact support, therefore,
by the Weierstrass Theorem $\mathcal{Q}(\pi)$ is bounded and attains its minimum.
Denote,
\begin{equation}
    \tau = \inf_{\|\pi\|=1} \mathcal{Q}(\pi).
    \label{pi*}
\end{equation}
In particular, for every $\pi$, $\mathcal{Q}(\pi)>0$ so we also have $\tau>0$. We fix from here on any constant
$$
C_\lambda \in \left(0,\ \underline{\beta}\,\tau\right).
$$
Consider the deterministic (conditional) covariance matrix
\begin{equation}
\mathscr{A}(x):=\mathbb{E}\!\left[\mathscr{A}_n(x)\,\middle|\,\mathscr{E}_I\right],
\end{equation}
so that for every unit vector $\pi$,
$$
\pi^\top\mathscr{A}(x)\pi=\mathbb{E}\!\left[\mathcal{Q}_n(\pi)\,\right]
\;\ge\;\underline{\beta}\,\mathcal{Q}(\pi)\;\ge\;\underline{\beta}\,\tau,
$$
where the first inequality is \eqref{bound of EQ_n}. Because $\mathscr{A}(x)$ is \emph{deterministic}, taking the infimum over $\|\pi\|=1$ is legitimate and yields
\begin{equation}
\lambda_{\min}\!\big(\mathscr{A}(x)\big)
=\inf_{\|\pi\|=1}\pi^\top\mathscr{A}(x)\pi
\;\ge\;\underline{\beta}\,\tau\;>\;C_\lambda\;>\;0 .
\label{eq:mean-pd}
\end{equation}
Similarly to Equation \eqref{eq:bernstein inequality}, by the scalar Bernstein inequality, for every $\varepsilon>0$ and every $\alpha,\beta$,
\begin{equation}
P\!\left(\big|\mathscr{A}_{n\alpha\beta}(x)-\mathscr{A}_{\alpha\beta}(x)\big|\ge\varepsilon
\Big|\mathscr{E}_I\right)\le 2e^{-c'\varepsilon^2\log n}=\frac{2}{n^{\,c'\varepsilon^2}}.
\label{eq:entry-conc}
\end{equation}
For an $A\times A$ matrix $B$, $\|B\|_{\mathrm{op}}\le\|B\|_F\le A\max_{\alpha,\beta}|B_{\alpha\beta}|$.
A union bound over the (finitely many, $n$-independent) $A^2$ entries therefore gives, from
\eqref{eq:entry-conc},
\begin{equation}
    P\!\left(\big\|\mathscr{A}_n(x)-\mathscr{A}(x)\big\|_{\mathrm{op}}\le A\varepsilon
\,\middle|\,\mathscr{E}_I\right)\ge 1-\frac{2A^2}{n^{\,c'\varepsilon^2}} .
\end{equation}
Note that $c'$ is the same constant as in \eqref{eq:bernstein inequality}: the constant is depending only on $\Phi$ and the support radius $s$, and the conditioning on $\mathscr{E}_I$ fixes the number of active samples $N_{B(\hat{x},sh_n)}=\Theta(\log n)$, so the Bernstein--Hoeffding exponent has the same dependence $c'\varepsilon^2\log n$ in both lemmas.

Choose $\varepsilon=\dfrac{\underline{\beta}\,\tau-C_\lambda}{A}$ (so $A\varepsilon=\underline{\beta}\,\tau-C_\lambda>0$) and set $\phi:=c'\varepsilon^2=\dfrac{c'\left(\underline{\beta}\,\tau-C_\lambda\right)^2}{A^2}>0$.
By Weyl's eigenvalue perturbation inequality for Hermitian matrices and \eqref{eq:mean-pd}, we obtain 
\begin{equation}
    \lambda_{\min}(\mathscr{A}_n(x))\;\ge\;\lambda_{\min}\!\big(\mathscr{A}(x)\big)
-\big\|\mathscr{A}_n(x)-\mathscr{A}(x)\big\|_{\mathrm{op}}
\;\ge\;\underline{\beta}\,\tau-\left(\underline{\beta}\,\tau-C_\lambda\right)\;=\;C_\lambda .
\end{equation}
Hence
$$
P\!\left(\lambda_{\min}(\mathscr{A}_n(x))\ge C_\lambda \Big|\mathscr{E}_I\right)
\;\ge\;1-\frac{2A^2}{n^{\phi}} ,
$$
which is the assertion (with the $n$-independent constant $2A^2$ in place of $2$). In
particular $\mathscr{A}_n(x)$ is positive definite and full rank for all sufficiently large $n$.
\end{proof}

\end{document}
